\chapter*{Abstract}
~\newline~\newline~
\markboth{Abstract}{Abstract}
\addcontentsline{toc}{chapter}{Abstract}

Solar coffee drying is a sustainable alternative for reducing the moisture content of coffee beans while decreasing dependence on conventional energy sources. However, its performance is strongly influenced by the variability of solar radiation, ambient temperature, relative humidity, wind conditions, and the thermal and geometric characteristics of the dryer. These factors may lead to non-uniform temperature and humidity distributions, prolonged drying times, and heterogeneous product quality. To address these challenges, this research proposes an integrated framework for the modeling, characterization, and intelligent control of thermo-hygrometric conditions in solar coffee dryers.\\

The drying system is modeled as an open, nonlinear, and multivariable dynamic system governed by the coupled mechanisms of solar radiation, heat conduction, convective heat transfer, moisture transport, and water evaporation. A multiscale modeling strategy is adopted by combining lumped-parameter physical models, Computational Fluid Dynamics (CFD) simulations, distributed experimental measurements, and data-driven models. High-fidelity CFD simulations are employed to characterize spatial distributions of temperature, airflow, and humidity, whereas reduced-order and hybrid models are intended to provide computationally efficient representations suitable for state estimation, prediction, and advanced control.\\

The proposed framework considers classical control, Model Predictive Control (MPC), and hybrid model-based control strategies. Unlike conventional approaches that regulate only average process variables, the proposed methodology also emphasizes the spatial uniformity of temperature and humidity as a performance criterion, while explicitly accounting for actuator constraints and time-varying environmental disturbances. Furthermore, the framework investigates the identification of sufficiently informative sensing and actuation architectures capable of reducing system complexity without compromising control performance.\\

The main contribution of this research is the integration of physical models, numerical simulations, experimental evidence, and data-driven techniques into a unified framework for thermo-hygrometric modeling and intelligent control. The proposed architecture explicitly distinguishes between experimentally validated operating conditions, simulation-supported operating regions, and extrapolation domains, providing a structured basis for assessing model reliability. It is expected that the proposed framework will contribute to the development of more uniform, adaptive, and computationally efficient solar coffee drying systems suitable for the environmental conditions of Panama.


\bigskip
 
\noindent\textit{Panamá, Agosto 2026}
\hfill A.~R.

